\documentclass[../../../include/open-logic-section]{subfiles}

\begin{document}

\olfileid{sth}{z}{sep}
\olsection{الفصل}

نبدأ بمبدأ يحل محل الاستيعاب الساذج:

\begin{axiom}[مخطط الفصل] لكل صيغة~$\phi(x)$، يكون الآتي مسلّمة:
لكل~$A$، توجد المجموعة~$\Setabs{x \in A}{\phi(x)}$.
\end{axiom}

لاحظ أن هذه ليست مسلّمة واحدة، بل \emph{مخطط} مسلّمات. فهناك
\emph{عدد لا نهائي} من مسلّمات الفصل، واحدة لكل صيغة~$\phi(x)$.
ويمكن كتابة المخطط على نحو مكافئ (وهو المعتاد) هكذا:

\begin{defish}
لكل صيغة~$\phi(x)$ لا تحتوي «$S$»، يكون الآتي مسلّمة:
\[
	\forall A \exists S \forall x(x \in S \liff (\phi(x) \land x \in A)).
\]
\end{defish}

وفقًا للاتفاق المذكور في بداية~\olref[sth][][]{part}، قد تحتوي
الصيغ~$\phi$ في مسلّمات الفصل على معاملات.\footnote{لشرح معنى ذلك،
  انظر المناقشة التي تلي مباشرة
  \olref[sfr][infinite][induction]{natinductionschema}.}

يبرر تصورنا التراكمي-التكراري للمجموعات الذي عرضناه مبدأ الفصل مباشرة.
ولرؤية السبب، لتكن~$A$ مجموعة. إذن تكونت~$A$ بحلول مرحلة ما~$S$
(بموجب~\stageshier). ولما كانت~$A$ قد تكونت في المرحلة~$S$، فقد
تكون جميع عناصر~$A$ قبل المرحلة~$S$ (بموجب~\stagesacc). وانظر الآن
خصوصًا في جميع المجموعات التي تنتمي إلى~$A$ وتحقق~$\phi$ أيضًا؛ فمن
الواضح أن جميع هذه المجموعات تكونت هي أيضًا قبل المرحلة~$S$. ولذلك
تُكوّن هي أيضًا في مجموعة~$\Setabs{x \in A}{\phi(x)}$ في المرحلة~$S$
(بموجب~\stagesacc).

وبخلاف الاستيعاب الساذج، يتفادى هذا مفارقة راسل. فلا يمكننا ببساطة
إثبات وجود المجموعة~$\Setabs{x}{x \notin x}$. بل إذا \emph{أُعطينا}
مجموعة~$A$، أمكننا إثبات وجود المجموعة
$R_A = \Setabs{x \in A}{x \notin x}$. لكن كل ما يثبته ذلك هو
$R_A \notin R_A$ و$R_A \notin A$، ولا يبعث أي منهما على قلق كبير.

غير أن للفصل نتيجة مباشرة ولافتة:

\begin{thm}\ollabel{thm:NoUniversalSet}
لا توجد مجموعة \emph{شاملة}، أي إن~$\Setabs{x}{x = x}$ لا توجد.
\end{thm}

\begin{proof}
برهانًا بالخلف، افترض أن~$V$ مجموعة شاملة. عندئذ توجد، بموجب الفصل،
$R = \Setabs{x \in V}{x \notin x} = \Setabs{x}{x \notin x}$،
وهذا يناقض مفارقة راسل.
\end{proof}

إن غياب مجموعة شاملة---بل انفتاح هرم المجموعات---من أكثر الأفكار
أساسية وراء التصور التراكمي-التكراري. ولذلك يجدر بنا أن نرى، حدسيًا،
أننا نستطيع بلوغ هذه النتيجة بطريق آخر. يجب أن تكون المجموعة الشاملة
!!a{element} في نفسها. لكن بحسب تصورنا التراكمي-التكراري، تظهر كل
مجموعة (للمرة الأولى) في الهرم عند أول مرحلة تأتي مباشرة بعد جميع
!!{element} التي تنتمي إليها. وينتج من ذلك أن \emph{لا} مجموعة تنتمي
إلى نفسها. فأي مجموعة تنتمي إلى نفسها يجب أن تظهر أول مرة مباشرة بعد
المرحلة التي ظهرت فيها أول مرة، وهذا محال. (سنرى في
\olref[sth][spine][rank]{defnsetrank} كيف نجعل هذا الشرح أدق باستعمال
مفهوم «رتبة» المجموعة. لكننا سنحتاج إلى وضع بضع مسلّمات أخرى قبل ذلك.)

وفيما يأتي بضع نتائج أخرى للفصل والامتدادية.

\begin{prop}\ollabel{prop:emptyexists}
إذا وجدت أي مجموعة، وجدت~$\emptyset$.
\end{prop}

\begin{proof}
إذا كانت~$A$ مجموعة، فإن
$\emptyset = \Setabs{x \in A}{x \neq x}$ توجد بموجب الفصل.
\end{proof}

\begin{prop}
توجد~$A \setminus B$ لكل مجموعتين~$A$ و$B$.
\end{prop}

\begin{proof}
توجد~$A \setminus B = \Setabs{x \in A}{x \notin B}$ بموجب الفصل.
\end{proof}

ويتضح أيضًا أن التقاطعات الاعتباطية (تقريبًا) موجودة:

\begin{prop}\ollabel{prop:intersectionsexist}
إذا كان~$A \neq \emptyset$، فإن
$\bigcap A = \Setabs{x}{(\forall y \in A)x \in y}$ توجد.
\end{prop}

\begin{proof}
لتكن~$A \neq \emptyset$، فيوجد~$c \in A$. عندئذ
$\bigcap A = \Setabs{x}{(\forall y \in A)x \in y} =
\Setabs{x \in c}{(\forall y \in A)x \in y}$، وهي موجودة بموجب الفصل.
\end{proof}

لكن لاحظ الشرط~$A \neq \emptyset$؛ إذ ستكون~$\bigcap \emptyset$
المجموعة الشاملة بصدق فارغ، وهذا يناقض
\olref{thm:NoUniversalSet}.

\end{document}
